% GENERATED FILE. Edit canonical lessons, then run npm run generate:books.
% canonical-source-hash: fnv1a64:54bec0291e88b117
% canonical-lessons: MW-C131-khichro, MW-C131-ker, MW-C131-sagri, MW-C131-gatta, MW-C131-papar

\chapter{Millet-and-lentil dish, Desert berry, Desert beans, Gram-flour dumplings, Papad}
\label{ch:mw-millet-and-lentil-dish-desert-berry-desert-beans-gram-flour-dumplings-papad}
\hlchaptermodality{131}

\begin{chapteropening}
\textbf{By the end of this chapter:} \emph{I can say a millet-and-lentil dish, a desert berry, desert beans, gram-flour dumplings and a papad.}

Complete the last lesson of chapter 131: I can say a millet-and-lentil dish, a desert berry, desert beans, gram-flour dumplings and a papad.
\end{chapteropening}

\section[\texorpdfstring{\mw{खीचड़ो} (khīchṛo)}{खीचड़ो (khīchṛo)}]{\mw{खीचड़ो} (khīchṛo) — a millet-and-lentil dish}

\label{lesson:MW-C131-khichro}



\begin{quote}
Before the new one: say the Marwadi for millet, then the Marwadi for a millet porridge.
\end{quote}

\subsection*{You'll want to know: \mw{खीचड़ो}}
\textbf{\mw{खीचड़ो}} — \emph{khīchṛo} — \textquotedblleft{}a millet-and-lentil dish\textquotedblright{}.

\begin{grammarlens}[title={What you've built}]
One more word for this chapter.
\end{grammarlens}

\subsection*{Your turn}
\begin{itemize}
\raggedright
  \item \emph{Say it:} \emph{khīchṛo}
  \item \emph{Say it:} \emph{khīchṛo}, once more
  \item \emph{Say it:} say \emph{khīchṛo}
  \item \emph{Recall:} say the Marwadi for millet, then the Marwadi for a millet porridge, then say \emph{khīchṛo} again
\end{itemize}

\begin{tcolorbox}[breakable,colback=teal!4,colframe=teal!35!black,title={Before you move on}]
What does \mw{खीचड़ो} mean? (\textquotedblleft{}A millet-and-lentil dish\textquotedblright{}.) Say it once more.
\end{tcolorbox}
\section[\texorpdfstring{\mw{केर} (ker)}{केर (ker)}]{\mw{केर} (ker) — a desert berry}

\label{lesson:MW-C131-ker}



\begin{quote}
Before the new one: say the Marwadi for a millet porridge, then the Marwadi for a millet-and-lentil dish.
\end{quote}

\subsection*{You'll want to know: \mw{केर}}
\textbf{\mw{केर}} — \emph{ker} — \textquotedblleft{}a desert berry\textquotedblright{}.

\begin{grammarlens}[title={What you've built}]
One more word for this chapter.
\end{grammarlens}

\subsection*{Your turn}
\begin{itemize}
\raggedright
  \item \emph{Say it:} \emph{ker}
  \item \emph{Say it:} \emph{ker}, once more
  \item \emph{Say it:} say \emph{ker}
  \item \emph{Recall:} say the Marwadi for a millet porridge, then the Marwadi for a millet-and-lentil dish, then say \emph{ker} again
\end{itemize}

\begin{tcolorbox}[breakable,colback=teal!4,colframe=teal!35!black,title={Before you move on}]
What does \mw{केर} mean? (\textquotedblleft{}A desert berry\textquotedblright{}.) Say it once more.
\end{tcolorbox}
\section[\texorpdfstring{\mw{सांगरी} (sā̃grī)}{सांगरी (sā̃grī)}]{\mw{सांगरी} (sā̃grī) — desert beans}

\label{lesson:MW-C131-sagri}



\begin{quote}
Before the new one: say the Marwadi for a millet-and-lentil dish, then the Marwadi for a desert berry.
\end{quote}

\subsection*{You'll want to know: \mw{सांगरी}}
\textbf{\mw{सांगरी}} — \emph{sā̃grī} — \textquotedblleft{}desert beans\textquotedblright{}.

\begin{grammarlens}[title={What you've built}]
One more word for this chapter.
\end{grammarlens}

\subsection*{Your turn}
\begin{itemize}
\raggedright
  \item \emph{Say it:} \emph{sā̃grī}
  \item \emph{Say it:} \emph{sā̃grī}, once more
  \item \emph{Say it:} say \emph{sā̃grī}
  \item \emph{Recall:} say the Marwadi for a millet-and-lentil dish, then the Marwadi for a desert berry, then say \emph{sā̃grī} again
\end{itemize}

\begin{tcolorbox}[breakable,colback=teal!4,colframe=teal!35!black,title={Before you move on}]
What does \mw{सांगरी} mean? (\textquotedblleft{}Desert beans\textquotedblright{}.) Say it once more.
\end{tcolorbox}
\section[\texorpdfstring{\mw{गट्टा} (gaṭṭā)}{गट्टा (gaṭṭā)}]{\mw{गट्टा} (gaṭṭā) — gram-flour dumplings}

\label{lesson:MW-C131-gatta}



\begin{quote}
Before the new one: say the Marwadi for a desert berry, then the Marwadi for desert beans.
\end{quote}

\subsection*{You'll want to know: \mw{गट्टा}}
\textbf{\mw{गट्टा}} — \emph{gaṭṭā} — \textquotedblleft{}gram-flour dumplings\textquotedblright{}.

\begin{grammarlens}[title={What you've built}]
One more word for this chapter.
\end{grammarlens}

\subsection*{Your turn}
\begin{itemize}
\raggedright
  \item \emph{Say it:} \emph{gaṭṭā}
  \item \emph{Say it:} \emph{gaṭṭā}, once more
  \item \emph{Say it:} say \emph{gaṭṭā}
  \item \emph{Recall:} say the Marwadi for a desert berry, then the Marwadi for desert beans, then say \emph{gaṭṭā} again
\end{itemize}

\begin{tcolorbox}[breakable,colback=teal!4,colframe=teal!35!black,title={Before you move on}]
What does \mw{गट्टा} mean? (\textquotedblleft{}Gram-flour dumplings\textquotedblright{}.) Say it once more.
\end{tcolorbox}
\section[\texorpdfstring{\mw{पापड़} (pāpaṛ)}{पापड़ (pāpaṛ)}]{\mw{पापड़} (pāpaṛ) — a papad}

\label{lesson:MW-C131-papar}



\begin{quote}
Before the new one: say the Marwadi for desert beans, then the Marwadi for gram-flour dumplings.
\end{quote}

\subsection*{You'll want to know: \mw{पापड़}}
\textbf{\mw{पापड़}} — \emph{pāpaṛ} — \textquotedblleft{}a papad\textquotedblright{}.

\begin{grammarlens}[title={What you've built}]
One more word for this chapter.
\end{grammarlens}

\subsection*{Your turn}
\begin{itemize}
\raggedright
  \item \emph{Say it:} \emph{pāpaṛ}
  \item \emph{Say it:} \emph{pāpaṛ}, once more
  \item \emph{Say it:} say \emph{pāpaṛ}
  \item \emph{Recall:} say the Marwadi for desert beans, then the Marwadi for gram-flour dumplings, then say \emph{pāpaṛ} again
\end{itemize}

\begin{tcolorbox}[breakable,colback=teal!4,colframe=teal!35!black,title={Before you move on}]
What does \mw{पापड़} mean? (\textquotedblleft{}A papad\textquotedblright{}.) Say it once more.
\end{tcolorbox}
